% Machine-assisted raw transcription, NOT visually verified formula by formula.
% Jethro van Ekeren, Homological Algebra, section 6.2, pages 119--130.
% This preserves extracted wording but formulas, diagrams, and some math symbols
% are imperfectly extracted. Do not mistake this for a certified verbatim edition.
% Requires \usepackage{fancyvrb} in the LaTeX preamble.
\section{\v{C}ech cohomology --- Jethro's notes (source transcription)}
\label{section-cech-jethro-transcription}

\subsection*{Source page 118}
\begin{Verbatim}[fontsize=\small]
118                        CHAPTER 6. COHOMOLOGY OF SHEAVES

                                                  · · ·


                      H2(X, F)          0         H2(X, G)


                      H1(X, F)          0         H1(X, G)


                  0      H0(X, F)      H0(X, I)      H0(X, G)


The middle column vanishes because I is injective.

There are two things we need to deduce from F being flabby (Exercise – for a reference, see
[15, Lemma 20.12.3]):


   1. That G is flabby,

   2. That the morphism Γ(X, I) →Γ(X, G) is surjective.

The second of these facts implies H1(X, F) = 0. So we conclude that H1(X, −) vanishes for
all flabby sheaves; this will be the base of an induction.

But the LES gives us isomorphisms Hk(X, G) →Hk+1(X, F) for all k ≥1, so since G is
flabby (first point above) we have Hk(X, F) = 0 for all k ≥1 by induction.

The proof of the third point is very similar to that of the second point.


6.2  Cech cohomology


Cech cohomology is a very convenient tool for computing sheaf cohomology. It’s what allows
you to turn sheaves on a space into concrete invariants like Betti numbers, etc. Of course,
not every open cover gives you the right answer (indeed {X} is an open cover of X, and
this would usually not be expected to work). In general one imagines that some cover by
“sufficiently small” open sets will give the right answer. A good reference is Hartshorne [6,
Section III.4].


6.2.1  Definition and elementary properties

Lemma 6.2.1. The restriction I|U of an injective sheaf I on X to an open subset U ⊂I, is
an injective sheaf on U.


Proof. Exercise.
\end{Verbatim}

\subsection*{Source page 119}
\begin{Verbatim}[fontsize=\small]
6.2. CECH COHOMOLOGY                                                   119

Definition 6.2.1. An open cover U = {Ui}i∈I of a topological space X is said to be locally
finite if, for every x ∈X, there exists an open V containing x such that only finitely many
sets V ∩Ui are non-empty.

Definition 6.2.2. A cover U of a topological space is said to be ’good’ if every finite inter-
section of opens from U is either empty or is connected and contractible to a point.

Let U = {Ui}i∈I be an open cover of a topological space X, and let F be a sheaf on X. For
any finite subset S = {i0, . . . , ip} ⊂I denote by US the open set ∩s∈SUs. Now we define the
abelian group
                            Cp(U, F) = Y  F(US).

                                        #S=p+1

We define Cech differentials
                            Cp(U, F) →Cp+1(U, F),
but to do so we have to fix an order on I. An element α ∈Cp(U, F) consists of an element

                                           αi0,...,ip ∈F(Ui0 ∩· · · ∩Uip) = F(US)
for each tuple i0 < . . . < ip, and dα ∈Cp+1(U, F) will consist of similar elements

                                         p+1
                                   (dα)i0,...,ip+1 = X (−1)j Res αi0,...,bij...,ip+1
                                           j=0

Here Res stands for the restriction from Ui0 ∩. . . ∩bUij ∩. . . ∩Uip+1 to the smaller open set
Ui0 ∩· · · ∩Uip+1. Also, although I suppose it’s obvious, let me emphasise: the sum here is
not happening inside any particular space of sections, the different Res αi0,...,bij...,ip+1 are just
entries in a big direct sum of spaces of sections.

Question: where, if anywhere, is local finiteness of U important?

DO AN EXAMPLE, SUCH AS THE CONSTANT SHEAF Z on S1 or on the torus.
Remark 6.2.2. We can describe Cp(U, F) as the global sections of a sheaf. For each inter-
section US denote jS : US →X the inclusion. Then we set
                             Cp(U, F) = Y  jS∗(F|US).
                                       #S=p+1

The Cech differentials are defined at this level
                                     ˇCp(U, F) →ˇCp+1(U, F),

and we obtain the standard Cech differentials by applying global sections term-by-term.
This looks similar to the general derived functor yoga: We have replaced our sheaf F by a
resolution

                          F


                             0      ˇC0      ˇC1    · · ·
\end{Verbatim}

\subsection*{Source page 120}
\begin{Verbatim}[fontsize=\small]
120                        CHAPTER 6. COHOMOLOGY OF SHEAVES


and then hit the resolution, term-by-term with our functor, in this case Γ(X, −), obtaining

                            0     ˇC0     ˇC1    · · ·

which we then take the cohomology. What is not clear is whether this correctly computes
the derived functors of Γ(X, −).

Lemma 6.2.3. Let U be any open cover of X. The complex of sheaves ˇC• (as defined above)
is a resolution of the sheaf F.


This lemma is [6, Chapter III, Lemma 4.2].


Proof. The first idea is that it suffices to check on stalks. Then one constructs a homotopy
from the identity to zero.


Actually this remark is useful, because it allows us to prove:

Lemma 6.2.4. Let U be any cover of X and F a flabby sheaf on X. Then

                  ˇH0(U, F) = Γ(X, F)  and  ˇHk(U, F) = 0  for k ≥1.             (6.2.1)


This lemma is [6, Chapter III, Lemma 4.3].

Proof. We use the resolution 0 →F →ˇC•.
Each ˇCp is defined as a product of pushforwards of restrictions of F. But products, pushfor-
wards, and restrictions of flabby sheaves are flabby. So every ˇCp is flabby.
This means ˇC• is a suitable complex of sheaves to compute Hk(X, F), just apply Γ(X, −)
term-by-term. Of course this recovers precisely the Cech complex with respect to the cover
U, so on the one hand we have ˇHk(X, F). On the other hand, what this is computing is
Hk(X, F), which is zero because F is flabby by hypothesis. So we have proved the result.



Recall that a flabby sheaf is acyclic. So Lemma 6.2.4 says that Cech cohomology gets the
right answer for flabby sheaves, for any choice of cover U (even the silly cover U = {X}).

We remark that products, pushforwards, and restrictions of injective sheaves are injective.
See [11, Proposition 2.4.1]1. Similarly for soft sheaves, I imagine the same proof yields the
following variant also.

Lemma 6.2.5. Let U be any cover of X and F a soft sheaf on X. Then

                  ˇH0(U, F) = Γ(X, F)  and  ˇHk(U, F) = 0  for k ≥1.             (6.2.2)

   1In these online notes it is claimed (proof of Proposition 5.10) that pushforwards of injective sheaves are
not necessarily injective
\end{Verbatim}

\subsection*{Source page 121}
\begin{Verbatim}[fontsize=\small]
6.2. CECH COHOMOLOGY                                                   121


Proof. The corresponding statement for fine sheaves is [1, Lemma 4.3]. TO DO: Look for
this and other statements in [8, Appendix A.7] or elsewhere.

From Wikipedia:

”For the Zariski topology on a Noetherian separated scheme, Cech and sheaf cohomology
agree for any quasi-coherent sheaf.”

Lemma 6.2.6 ([6, III, Lemma 4.4]). Let X be a topological space, U a cover of X and F a
sheaf on X. There is a canonical morphism
                                  ˇHp(U, F) →Hp(X, F)
for each p ∈Z+.

Proof. The idea is grab a resolution I• of F by injective sheaves, as well as the Cech resolution
ˇC• (as sheaves) from above. By the magic of injective objects there is a morphism ˇC• →I•,
and furthermore by the usual chain homotopy argument, the induced map in cohomologies
is always the same.

The big question is whether Cech cohomology computes sheaf cohomology correctly for
sheaves in general, not just flabby ones. This will depend in general on how good a cover U
we take. The right notion, it turns out, is the following definition.

Definition 6.2.3. Let F ∈Sh(X) and let U = {Ui}i∈I be an open cover of X. We say U is
F-acyclic if
                 Hk(US, F) = 0   for all k ≥1 and all finite S ⊂I.

Let us remark that, while ”acyclic sheaves” and ”F-acyclic covers” are related, one should
be careful to keep straight the relationship between the notions. Saying ”F is acyclic” is the
same as saying ”the cover {X} is F-acyclic”. On the other hand, saying ”U is acyclic for
F” is the same as saying that every restriction F|US is acyclic (as a sheaf on US). A natural
question is whether ”F is acyclic” implies ”every cover U of X is F-acyclic”. The answer is
no, and it is easy to see why: consider the constant sheaf Z on X = R2, it is acyclic. Now
consider a cover containing an open annulus, then the sheaf cohomology of the constant
sheaf on the annulus is nontrivial, so the cover is not Z-acyclic.


6.2.2  The Leray theorem

Above I said that the ’right notion’ of cover appropriate to compute sheaf cohomology of F
was in Definition 6.2.3. This is the content of Leray’s theorem.

Theorem 6.2.7 (Leray’s theorem). Let F be a sheaf on X and suppose U is an open cover
of X which is F-acyclic. Then the natural morphism
                                 ˇHk(U, F) →Hk(X, F)

is an isomorphism for all k ≥0.
\end{Verbatim}

\subsection*{Source page 122}
\begin{Verbatim}[fontsize=\small]
122                        CHAPTER 6. COHOMOLOGY OF SHEAVES

Proof. We choose a resolution of F by injective sheaves In. We set up a bicomplex

                             Kp,q = ˇCp(Iq) = Y  Iq(US).

                                          #S=p+1

                                                                                                                            ...                                                               ...                                                               ...


                         ˇC0(I2)        ˇC1(I2)        ˇC2(I2)            · · ·


                         ˇC0(I1)        ˇC1(I1)        ˇC2(I1)            · · ·


                         ˇC0(I0)        ˇC1(I0)        ˇC2(I0)            · · ·


This will be page E0 of a spectral sequence, in which we produce page E1 by taking vertical
cohomology first.  In column p, every Ep,q0  has a summand corresponding to US and the
vertical differential respects this splitting into summands. So

                                  Ep,q1 = Y  Hq(I•(US)).
                                    #S=p+1

The sheaves Iq|US are injective by Lemma 6.2.1, so by definition of derived functors, we have

                                  Ep,q1 = Y  Hq(US, F).
                                    #S=p+1

Since U is F-acyclic by hypothesis, we have Ep,q1 = 0 for q ≥1, and Ep,01 = Q #S=p+1 F(US).
Evidently this is just ˇCp(F) The morphisms of the E1-page are the Cech differentials, too.
So Ep,02 = ˇHp(X, F) and the spectral sequence collapses here.

Now we consider the spectral sequence with the same E0-page, but first we take cohomology
with respect to the horizontal differentials. Since Iq is injective it is flabby, and so by Lemma
6.2.4 the complex

          0      Γ(X, Iq)        ˇC0(Iq)        ˇC1(Iq)        ˇC2(Iq)            · · ·

is exact (i.e., ˇHp(U, Iq) = 0 for k ≥1). Thus the E1-page has E0,q1 = Γ(X, Iq), with vertical
maps being the maps Iq →Iq+1 at the level of global sections. Since the Iq are injective, this
complex computes Hq(X, F) essentially by definition. In other words E0,q2 = Hq(X, F).
So we get the desired conclusion Hk(X, F) ∼= ˇHk(U, F).
\end{Verbatim}

\subsection*{Source page 123}
\begin{Verbatim}[fontsize=\small]
6.2. CECH COHOMOLOGY                                                   123


It is possible to talk about ’the’ Cech cohomology of a sheaf on a topological space, without
reference to a choice of cover. If U′ is a refinement of U then there exist induced maps

                           Hk(U, F) →Hk(U′, F).

The collection of all (locally finite) covers forms a category, and we define

Definition 6.2.4. The Cech cohomology of F on X (not relative to any particular choice
of cover), is the colimit
                        Hk(X, F) = colim Hk(U, F).
                                          U

Then of course any condition (such as Theorem 6.2.7) which guarantees an isomorphism
ˇHk(U, F) →Hk(X, F) for SOME cover U will guarantee an isomorphism ˇHk(X, F) →
Hk(X, F), because the colimit gets sandwiched in between.

So suppose we have a sheaf F and we want to compute its sheaf cohomology. What do we
do? The most obvious thing is to seek a cover U of X which is F-acyclic.  If U consists
of finitely many sets and we can determine the F(Ui) explicitly, then in principle we can
compute the Cech cohomology, and we know it coincides with the sheaf cohomology. This
turns out to be a good way to tackle sheaf cohomology of constant sheaves in the smooth
topology. We have the following result.
Proposition 6.2.8. Let U be a good open cover of X, then U is F-acyclic for F = Z the
constant sheaf.

Proof. FIND A PROOF.

On the other hand an algebraic variety with the Zariski topology typically has no good
covers at all. And more generally, for many standard sheaves of interest F there is no useful
F-acyclic cover.

The answer in these cases, is to resolve our sheaf F by a complex of sheaves Fi, and use a
cover which is Fi-acyclic for each i. There is some homological algebra machinery which then
recovers the Hk(X, F) in terms of the Fi (BE MORE PRECISE). In practice, paradoxically,
on usually goes to an extreme: resolve F in a class of sheaves with respect to which the cover
{X} is acyclic, i.e., we resolve F by acyclic sheaves Fi.

In the case of a smooth manifold, opens homeomorphic to balls in Rn usually suffice, for qua-
sicoherent sheaves in algebraic geometry, a cover for which all opens and their intersections
are open affines in X is sufficient [15, Lemma 30.2.6].


6.2.3  Example: Sheaf cohomology on the elliptic curve

Now let’s do some examples. Let E = C/Z + Zτ be a smooth elliptic curve. Let’s compute
the Cech cohomology of OE using the cover U1 = E\p1 and U2 = E\p2. The Cech complex
has only two terms
                               C0 →C1
\end{Verbatim}

\subsection*{Source page 124}
\begin{Verbatim}[fontsize=\small]
124                        CHAPTER 6. COHOMOLOGY OF SHEAVES

where C0 = OE and
                         C0 = j1∗OE|U1 ⊕j2∗OE|U2,

and
                            C1 = j1,2∗OE|U1∩U2.

We can describe H0(U1, OE) in terms of elliptic functions. Write Λ = Z + Zτ, and recall the
Weierstrass function
                              1              1                        ℘(z) =  + X      −1     ,
                                  z2            (z −λ)2   λ2
                                          λ∈Λ\{0}

which satisfies
                        ℘(z + 1) = ℘(z),    ℘(z + τ) = ℘(z),

i.e., is an elliptic function.  It turns out there is no elliptic function with a simple pole at
z = 0, so ℘(z) together with all its derivatives (and the constant function) must form a basis
of H0(E\{0}, OE). Hence H0(U1, OE) is spanned by 1, ℘(z −p1) and all its derivatives.

The space H0(U1 ∩U2, OE) contains the restrictions of all the functions on U1 and U2 already
seen, and essentially one more. Consider

                             1            1     1    z
                         ζ(z) =  + X      +  +        ,
                                z           z −λ  λ   λ2
                                         λ∈Λ\{0}

which is well-defined and meromorphic on C, but satisfies

                      ζ(z + 1) = ζ(z) + α,     ζ(z + τ) = ζ(z) + β,

for some constants α and β. Thus ζ(z) is not elliptic but

                       ζp1,p2(z) = ζ(z −p1) −ζ(z −p2) ∈H0(U1 ∩U2, OE)

is. By the residue theorem, any elliptic function with poles at p1 and p2 is a linear combination
of ζp1,p2, the constant function 1, and derivatives of ℘(z−p1) and ℘(z−p2). The Cech complex
is
                    H0(U1, OE) ⊕H0(U2, OE) →H0(U1 ∩U2, OE)
where (f, g) 7→f|U1∩U2 + g|U1∩U2. Clearly the kernel is C(1, −1) and the cokernel is Cζp1,p2.
Thus we have proved H0(OX) ∼= H1(OX) ∼= C.

Now in
                             Γ(E, C0) →Γ(E, C1)

the kernel of the differential consists of multiples of (1, 1) (the constant function in both
entries). The cokernel consists of multiples of ζ(z −p1) −ζ(z −p2).

We have just computed that

                         (C  n = 0, 1
                         RnΓ(E, OE) ∼=
                                          0    i ≥2.
\end{Verbatim}

\subsection*{Source page 125}
\begin{Verbatim}[fontsize=\small]
6.2. CECH COHOMOLOGY                                                   125

6.2.4  Fine sheaves and the de Rham theorem

Definition 6.2.5. Let F be a sheaf over a topological space X. We say F is a fine sheaf if,
for every open U ⊂X and every locally finite open cover {Vi}i∈I of U, there exist morphisms

                                    φi : F(Vi) →F(U)

such that
          X φi ◦ρUVi = IdF(U) .
                                                                 i

This definition makes sense over a general topological space, but the most common context
in which fine sheaves can be constructed is that of smooth real manifolds, using ’partitions
of unity’.

Definition 6.2.6. Let M be a smooth manifold and {Ui}i∈I an open cover of M. A partition
of unity, subordinate to {Ui}i∈I consists of:

   1. for each i ∈I, a smooth function βi : M →R whose support is contained inside some
     compact subset Ki of Ui,


such that
            X φi = 1.

                                                   i∈I

Suppose F is a sheaf of C∞(M)-modules, then a choice of partition of unity on each U ⊂X
gives us a collection of φi : F(Vi) →F(U) as in the definition of fine sheaf, i.e., F is a fine
sheaf.

Indeed let f ∈F(Vi) and consider the product βi · f, whose support is contained in Ki. We
take any cover of U\Ki (which is open), together with Vi, to get a cover of U. Consider
sections βi · f on Vi and 0 on all the other opens, which by the sheaf axioms glue to give a
section in F(U) which we denote φi(f). This gives our morphism φi : F(Vi) →F(U).

So we have shown

Proposition 6.2.9. A sheaf of C∞(M)-modules over a smooth manifold M is a fine sheaf.


Proof. Well, it follows from the existence of partitions of unity on smooth manifolds. See
any textbook on differential geometry.


A key property of fine sheaves is the following.

Proposition 6.2.10. Let F be a fine sheaf on a topological space X. Let U = {Ui}i∈I be
any locally finite open cover of X. Then the Cech cohomology satisfies

                              ˇHk(U, F) = 0  for all k ≥1.
\end{Verbatim}

\subsection*{Source page 126}
\begin{Verbatim}[fontsize=\small]
126                        CHAPTER 6. COHOMOLOGY OF SHEAVES


Proof.

In other words if F is a fine sheaf, then every cover U of X is F-acyclic. In particular, taking
U = {X}, we conclude that any fine sheaf is acyclic.

Let F be a sheaf, and let us consider a resolution F →F•.  In particular the augmented
complex

                  0     F      F0      F1      F2            · · ·

is exact. Suppose there exists a cover U of X such that on every intersection US the corre-
sponding complex of sections is exact, i.e.,

       0      F(US)      F0(US)      F1(US)      F2(US)            · · ·         (6.2.3)

Notice that exactness of the sequence of sheaves says that for every j ≥0 and for every
point x ∈X there is a neighbourhood U of x such that Fj−1(U) →Fj(U) →Fj+1(U) is
exact. The condition stated above is sort of saying that there be a particular cover U that
implements exactness ’uniformly’ in j and x.

Theorem 6.2.11. Let F be a sheaf on X, and 0 →F →F• a resolution. Let U be a cover
of X. We assume

   1. the cover U is Fi-acyclic for each i,

   2. the exactness of (6.2.3) holds for U.

Then there exist isomorphisms
                               ˇHk(U, F) ∼= Hk(Γ(X, F•)).                            (6.2.4)

Proof. We set up a bicomplex

                             Kp,q = ˇCp(Fq) = Y  Fq(US).

                                          #S=p+1

                                                                                                                          ...                                                                  ...                                                                  ...


                       ˇC0(F2)       ˇC1(F2)       ˇC2(F2)            · · ·


                       ˇC0(F1)       ˇC1(F1)       ˇC2(F1)            · · ·


                       ˇC0(F0)       ˇC1(F0)       ˇC2(F0)            · · ·
\end{Verbatim}

\subsection*{Source page 127}
\begin{Verbatim}[fontsize=\small]
6.2. CECH COHOMOLOGY                                                   127


This will be page E0 of a spectral sequence, in which we produce page E1 by taking vertical
cohomology first.  In column p, every Ep,q0  has a summand corresponding to US and the
vertical differential respects this splitting into summands. So

                                  Ep,q1 = Y  Hq(F•(US)).
                                   #S=p+1

Our assumption 6.2.3 is precisely that each of these summands vanishes for q ≥1. So we
have Ep,q1 = 0 for q ≥1.  Also Ep,01 = Q #S=p+1 F(US).  Evidently this is just ˇCp(U, F).
The morphisms of the E1-page are the Cech differentials, too. So Ep,02 = ˇHp(U, F) and the
spectral sequence collapses here.

Now we consider the spectral sequence with the same E0-page, but first we take cohomology
with respect to the horizontal differentials. Of course now Ep,q1 = ˇHp(U, Fq). By hypothesis
each Fq  is U-acyclic, so we have ˇHp(U, Fq) = 0 for all p ≥1.  Thus the E1-page has
E0,q1  = Γ(X, Fq), with vertical maps being the maps Fq →Fq+1 at the level of global
sections. Thus E0,q2 = hq(Γ(X, F•)).
So we get the desired conclusion ˇHk(U, F) ∼= Hk(Γ(X, F•)).


With this theorem it may seem that we still haven’t gotten out the the flybottle, because
neither side of the isomorphism is a sheaf cohomology. The Leray theorem says, however,
that the LHS coincides with sheaf cohomology if the cover U is F-acyclic. This also may
not seem like much progress, because it might be hard or impossible to find a cover which is
acyclic for our sheaf. Well, life is hard. But in fact we are in a better position now than we
were before.

What one tends to do is choose the Fi from a class of sheaves which are acyclic with respect
to any cover (so we take the Fi to be injective, flabby, soft or fine sheaves), so the first
hypothesis in the theorem is taken care of.

Corollary 6.2.12. Let F be a sheaf on X, and 0 →F →F• a resolution. Let U be a cover
of X. We assume


   1. the cover U is F-acyclic,

   2. the sheaves Fi are fine or soft or whatever (guaranteeing U is acyclic),

   3. the exactness of (6.2.3) holds for U.


Then there exist isomorphisms

                         Hk(X, F) ∼= Hk(Γ(X, F•)).                            (6.2.5)

For example we consider the constant sheaf R on a manifold. Any good cover is R-acyclic.
If we resolve R by C∞(M)-modules then, because these are fine, we get an isomorphism
between sheaf cohomology of R and cohomology of global sections of our resolution.
\end{Verbatim}

\subsection*{Source page 128}
\begin{Verbatim}[fontsize=\small]
128                        CHAPTER 6. COHOMOLOGY OF SHEAVES

6.2.5  The Poincar´e lemma


State Poincar´e lemma.

Use it to prove the de Rham isomorphism.

If U is a good cover of the smooth manifold, then U is ΩkM-acyclic for all k.

Proposition 6.2.13. Let U be a good open cover of X, then U is F-acyclic for F = R the
constant sheaf.


Proof. Basically IS the statement of the Poincar´e lemma.


6.2.6  Comparison with singular cohomology


Griffiths and Harris has a nice presentation, it seems. See [4, p. 42] (de Rham theorem on
next page too).

Move this example of nerve of a cover back to the algebraic topology section. Also think
about whether it’s better to work with ASC or simplicial set.

Let U = {Ui}i∈I be a good open cover of X. We can associate with it an ASC K based
on the index set V ert(K) = I, namely S ⊂I is in K if US = ∩i∈SUi is non empty. This is
called the nerve of the cover.

We then have the homology complexes Cori• (K) and Csing•  (|K|) which we know to be homotopy
equivalent.

Proposition 6.2.14. The dual complexes C•ori(K) and C•sing(|K|) are homotopy equivalent.


Proof. We recall that the dual complexes in each case, are the term-wise duals, with adjoint
maps as differentials. We leave the proof of homotopy equivalence as an exercise.

Now we argue that ˇCn(U, Z) coincides with Cnori(K) and in fact the complexes coincide. It’s
actually pretty easy to see if you think about it.

According to [7, Corollary 4G.3] if K is the nerve of a good cover of X then |K| is homotopy
equivalent to X.

PICTURE OF 6 HEMISPHERES COVERING THE SPHERE. AND OCTAHEDRON.

So this proves

Theorem 6.2.15. Let U be a good open cover of X, then

                                ˇHk(U, R) ∼= Hksing(X, R).
\end{Verbatim}

\subsection*{Source page 129}
\begin{Verbatim}[fontsize=\small]
6.2. CECH COHOMOLOGY                                                   129

6.2.7  Cohomology as a special case of derived pushforward

Derived functors of pushforward to a point have a very natural interpretation.  If X  is a
sufficiently nice topological space they recover singular cohomology of X.

Theorem 6.2.16. Let X be a topological manifold, then there are isomorphisms
                          RnΓ(X, RX) ∼= Hnsing(X).

The theorem is true for topological spaces much more general than topological manifolds. It
is true (for example) if X is paracompact and locally contractible.

Although it is important to know that enough projectives or enough injectives exist, in the
day-to-day business of computing derived functors it is often possible to get away with using
less exotic objects. Given a left exact functor F, there is a notion of F-injective object, and
it is a theorem that a resolution by F-injective objects suffices to compute derived functors
RiF.

In the case of sheaves, one has the following notion:

Definition 6.2.7. A sheaf F on X is said to be flabby if the restriction morphism F(X) →
F(U) is surjective for every open set U ⊂X.

It is then a standard theorem (I found it in these notes [17, Proposition 6.9], look up a
textbook reference)

Theorem 6.2.17. A flabby sheaf is Γ(X, −)-injective, indeed Γ(X, −)-acyclic (which is
stronger).

The sheaves of singular cochains are not necessarily injective, but  if X  is a topological
manifold then they are flabby. So by the homological algebra general nonsense above, these
sheaves can be used to resolve if one is computing global sections.

The proof of Theorem 6.2.16 then proceeds by introducing sheaves of singular cochains.
Roughly speaking: for n ∈Z≥0 a basic singular chain in U ⊂X is a continuous function
σ  : ∆n →U, where ∆n is the standard n-simplex, and a singular chain is a formal linear
combination of finitely many basic ones. The vector space of singular chains is denoted Cn(U)
and singular cochains are elements of Cn(U) = Cn(U)∗the dual vector space. These form
presheaves, which we then sheafify to obtain sheaves Cnsing. The usual coboundary produces
a complex
                          0 →C0sing →C1sing →C2sing →· · ·
of sheaves which resolves RX (this is proved by passing to small open sets in X which are
sufficiently nice to check explicitly). A key point is that the sheaves Cnsing are flabby, so
the derived global sections of RX are computed as the cohomology of the complex of vector
spaces
                  0 →Γ(X, C0sing) →Γ(X, C1sing) →Γ(X, C2sing) →· · ·
Now these vector spaces are exactly the vector spaces of singular cochains which are used to
define singular cohomology in basic topology.
\end{Verbatim}

\subsection*{Source page 130}
\begin{Verbatim}[fontsize=\small]
130                        CHAPTER 6. COHOMOLOGY OF SHEAVES

6.2.8  Schemes and the Zariski topology

The main theorem is [6, Chap. III, Theorem 4.5]

Theorem 6.2.18. Let X be a Noetherian separated scheme and U a cover of X by open
affine sets. Let F be a quasi-coherent sheaf on X. Then the natural morphisms

                                 ˇHk(U, F) →Hk(X, F)

are isomorphisms.


6.2.9 GAGA

So is the Zariski topology good enough or not? Well, yes and no.

Serre showed that if we are interested in coherent sheaves on complex algebraic varieties,
then sheaf cohomology in the Zariski and analytic topologies coincide.

But note that sheaf cohomology of constant sheaves in the Zariski topology are NOT well
behaved, even for complex algebraic varieties.

Etale cohomology was invented as a better topology to use in positive characteristic. I guess
I’m curious: does it also solve the issue about the constant sheaf, over complex algebraic
varieties?

Yeah, according to Kedlaya’s notes https://kskedlaya.org/18.726/etale-cohom.pdf over
the etale site (for a complex algebraic variety X) there are ℓ-adic sheaves which behave a lot
like the constant sheaf, and whose cohomology reproduces the Betti numbers of X
\end{Verbatim}